% Source: https://artofproblemsolving.com/community/c3h1898891_zeometry
% Stable identity: aops-topic-1898891-q01
Name the inscribed angles by the arcs they stand on:
\[
x=\angle CAD=\angle CBD,\quad \beta=\angle ACD,\quad \gamma=\angle BDC,\quad \delta=\angle ADB=\angle ACB .
\]

\begin{asy}
usepackage("amsmath");
size(6.5cm);
defaultpen(fontsize(10pt));
pair A=dir(110),B=dir(140),C=dir(200),D=dir(320),O=(0,0);
pair P=extension(D,A,C,B), Q=extension(A,C,B,D);
draw(unitcircle,gray(0.45));
draw(A--B--C--D--cycle,linewidth(0.9));
draw(A--C,linewidth(0.7));draw(B--D,linewidth(0.7));
draw(A--P,dashed);draw(B--P,dashed);
draw(C--P,blue+linewidth(1.1));draw(C--D,blue+linewidth(1.1));draw(D--Q,blue+linewidth(1.1));
void tk(pair a,pair b,int k){pair m=(a+b)/2,u=unit(b-a),n=(-u.y,u.x);
 for(int i=0;i<k;++i){pair c=m+(i-(k-1)/2)*0.05*u; draw((c-0.06*n)--(c+0.06*n),blue+linewidth(0.9));}}
tk(C,D,1);tk(C,P,1);tk(D,Q,1);
dot(A);dot(B);dot(C);dot(D);dot(P);dot(Q);
label("$A$",A,NE);label("$B$",B,dir(160));label("$C$",C,SW);label("$D$",D,SE);
label("$P$",P,NW);label("$Q$",Q,1.4*E);
draw(arc(A,0.22,degrees(C-A),degrees(D-A)),red+linewidth(0.7));
label("$x$",A+0.32*dir((degrees(C-A)+degrees(D-A))/2),red);
label("\mbox{}",(min(currentpicture,user=true)+max(currentpicture,user=true))/2,black); // keeps the text after the figure black when LaTeX typesets the labels inline
\end{asy}

\textbf{Triangle $ACD$.} Its angles are $x$, $\beta$ and $\angle ADC=\angle ADB+\angle BDC=\delta+\gamma$, so
\[
x+\beta+\gamma+\delta=180^\circ. \tag{1}
\]
\textbf{Triangle $DQC$.} Since $Q$ lies on the diagonals, $\angle QDC=\angle BDC=\gamma$ and $\angle DCQ=\angle DCA=\beta$. As $DQ=DC$, the angles opposite these equal sides are equal: $\angle DQC=\angle DCQ=\beta$. So
\[
\gamma+2\beta=180^\circ. \tag{2}
\]
\textbf{Triangle $PDC$.} Here $P$ lies on ray $DA$ beyond $A$ and on ray $CB$ beyond $B$, so $\angle PDC=\angle ADC=\delta+\gamma$ and $\angle PCD=\angle BCD=\angle BCA+\angle ACD=\delta+\beta$. As $CP=CD$, $\angle CPD=\angle CDP=\delta+\gamma$, and the angle sum gives
\[
2(\delta+\gamma)+(\delta+\beta)=180^\circ,\quad\text{i.e.}\quad 3\delta+2\gamma+\beta=180^\circ. \tag{3}
\]
From (2), $\gamma=180^\circ-2\beta$. Then (3) gives $3\delta+360^\circ-3\beta=180^\circ$, so $\delta=\beta-60^\circ$. Finally (1) gives
\[
x=180^\circ-\beta-(180^\circ-2\beta)-(\beta-60^\circ)=60^\circ .
\]
So $\angle CAD=\boxed{60^\circ}$. Such quadrilaterals exist: for example $\beta=75^\circ$, $\gamma=30^\circ$, $\delta=15^\circ$ (arcs $AB,BC,CD,DA$ of $30^\circ,60^\circ,120^\circ,150^\circ$), which is the configuration drawn.
